\documentclass{njtech-thesis}

\thesisyear{2026}
\thesismonth{6}
\thesistitle{生成式人工智能使用深度对大学生学习绩效的影响研究——基于批判性思维与AI依赖的视角}
\thesismajor{国际经济与贸易}
\thesisclass{国贸2201}
\thesisauthor{张三}
\thesisadvisor{李四\quad 教授}
\thesisdate{2025.12--2026.6}

\begin{document}

\makecover

\makedeclaration

\makefrontmatter

\begin{cnabstract}{随着生成式人工智能（Generative Artificial Intelligence，GAI）技术的快速发展，以ChatGPT、DeepSeek、文心一言等为代表的生成式人工智能工具正在深刻改变大学生的学习方式与知识获取模式。尤其在高等教育领域，GAI已广泛应用于资料搜集、语言翻译、论文写作、逻辑分析以及课程学习等多个场景。与此同时，生成式人工智能对大学生学业发展的影响也呈现出明显的"双刃剑"特征：一方面，GAI能够提升学习效率、优化知识理解并增强自主学习能力；另一方面，过度依赖AI可能导致认知卸载、思维惰性以及学术诚信风险。因此，深入探讨大学生生成式人工智能使用行为对学习绩效的影响机制，具有重要的理论价值与现实意义。基于技术接受理论、认知负荷理论与建构主义学习理论，本文构建了"AI使用深度—学习绩效"的分析框架，并引入批判性思维与AI依赖作为重要影响因素，重点考察不同AI使用方式对大学生学习效果的差异影响。本文将AI使用深度界定为学生在使用生成式人工智能过程中所表现出的认知参与程度、互动质量以及高阶学习投入水平，强调其与简单"使用频率"之间的本质区别。本文采用问卷调查法，以高校国际经济与贸易专业本科生为主要研究对象，通过Likert五级量表收集大学生在AI使用行为、批判性思维、AI依赖以及学习绩效等方面的数据，并运用SPSS与Stata软件进行描述性统计分析、信度效度检验以及多元回归分析。研究重点检验AI使用深度是否显著提升大学生学习绩效，以及批判性思维与AI依赖在其中所发挥的作用。研究认为，生成式人工智能对大学生学习的影响并非简单的"促进"或"抑制"，而取决于学生如何使用AI工具。深层次、互动式、反思型的AI使用能够有效促进知识理解与学习能力提升，而浅层、依赖型使用则可能削弱学生独立思考能力。基于研究结论，本文进一步提出高校应加强大学生AI素养教育，优化课程设计与教学评价体系，引导学生形成"人机协同"的学习模式，从而实现生成式人工智能与高等教育的深度融合。}
生成式人工智能；AI使用深度；学习绩效；批判性思维；AI依赖；大学生
\end{cnabstract}

\begin{enabstract}{With the rapid development of Generative Artificial Intelligence (GAI) technology, generative AI tools represented by ChatGPT, DeepSeek, and Wenxin Yiyan are profoundly transforming college students' learning methods and knowledge acquisition patterns. In the field of higher education, GAI has been widely applied in multiple scenarios such as data collection, language translation, paper writing, logical analysis, and course learning. Meanwhile, the impact of generative AI on college students' academic development exhibits a distinct "double-edged sword" characteristic: on one hand, GAI can enhance learning efficiency, optimize knowledge understanding, and strengthen autonomous learning ability; on the other hand, excessive reliance on AI may lead to cognitive offloading, thinking inertia, and academic integrity risks. Therefore, it is of significant theoretical value and practical significance to deeply explore the mechanism of college students' generative AI usage behavior on learning performance. Based on the Technology Acceptance Theory, Cognitive Load Theory, and Constructivist Learning Theory, this paper constructs an analytical framework of "AI Usage Depth — Learning Performance," and introduces critical thinking and AI dependency as important influencing factors, focusing on the differential effects of different AI usage patterns on college students' learning outcomes. This paper defines AI usage depth as the level of cognitive engagement, interaction quality, and higher-order learning investment that students exhibit in the process of using generative AI, emphasizing its essential distinction from simple "usage frequency." This paper adopts the questionnaire survey method, taking undergraduate students majoring in International Economics and Trade as the main research subjects, collecting data on college students' AI usage behavior, critical thinking, AI dependency, and learning performance through Likert five-point scales, and conducting descriptive statistical analysis, reliability and validity testing, and multiple regression analysis using SPSS and Stata software. The research focuses on examining whether AI usage depth significantly improves college students' learning performance, and the roles played by critical thinking and AI dependency. The study concludes that the impact of generative AI on college students' learning is not simply "promoting" or "inhibiting," but depends on how students use AI tools. Deep, interactive, and reflective AI usage can effectively promote knowledge understanding and learning ability improvement, while shallow, dependent usage may weaken students' independent thinking ability. Based on the research conclusions, this paper further proposes that universities should strengthen college students' AI literacy education, optimize curriculum design and teaching evaluation systems, and guide students to form a "human-machine collaboration" learning model, thereby achieving deep integration of generative AI and higher education.}
Generative Artificial Intelligence; AI Usage Depth; Learning Performance; Critical Thinking; AI Dependency; College Students
\end{enabstract}

\tableofcontents
\newpage

\makebody

\section{绪论}

\subsection{研究背景和意义}

随着以ChatGPT、DeepSeek、文心一言为代表的生成式人工智能（Generative Artificial Intelligence，GAI）技术的突破性进展，人工智能正在从传统的"辅助工具"向"智能伙伴"转变，深刻重塑人类的知识生产与获取方式。2022年11月OpenAI发布ChatGPT以来，生成式人工智能在全球范围内引发了前所未有的关注与应用浪潮。据联合国教科文组织2024年报告显示，全球已有超过50\%的高校学生在学习过程中使用过生成式AI工具。在中国，中国互联网络信息中心（CNNIC）发布的第53次《中国互联网络发展状况统计报告》显示，截至2024年6月，我国网民规模达10.99亿，其中使用AI工具的比例持续攀升，大学生群体是AI工具使用最为活跃的人群之一。

在国际经济与贸易专业领域，生成式人工智能的应用场景尤为广泛。学生可以利用GAI进行国际贸易术语查询、跨境市场分析、外文文献翻译、贸易政策解读以及商务谈判模拟等。然而，GAI在为学习带来便利的同时，也引发了一系列值得关注的问题：部分学生过度依赖AI完成作业和论文，导致独立分析能力下降；AI生成内容的准确性问题可能误导学生的学习方向；AI使用的边界模糊也给学术诚信带来了新的挑战。

\subsubsection{理论意义}

从理论层面看，现有关于技术与教育的研究多聚焦于AI的"使用与否"或"使用频率"，而对AI使用的"深度"与"质量"关注不足。本文将"AI使用深度"这一概念引入教育技术研究领域，从认知参与、互动质量和高阶学习投入三个维度进行界定与测量，有助于丰富技术接受理论与认知负荷理论在AI教育应用领域的理论内涵。同时，本文将批判性思维与AI依赖纳入分析框架，揭示二者在AI使用深度与学习绩效之间的调节作用机制，为理解"技术—认知—绩效"的复杂关系提供了新的理论视角。

\subsubsection{实践意义}

从实践层面看，本研究聚焦国际经济与贸易专业本科生的AI使用行为，研究结论可为高校制定AI使用规范、优化课程设计以及完善教学评价体系提供实证依据。同时，本研究提出的"人机协同"学习模式建议，有助于引导大学生形成健康的AI使用习惯，在享受技术红利的同时保持独立思考能力，对于推动高等教育数字化转型具有重要的参考价值。

\subsection{国内外研究现状}

国内外学者针对生成式人工智能在教育领域的应用已开展了较为丰富的研究，但关于AI使用深度与学习绩效之间关系的研究仍处于起步阶段。

\subsubsection{生成式人工智能教育应用研究}

生成式人工智能在教育领域的应用研究主要集中在以下几个方面：一是AI辅助教学的效果评估。Kasneci等（2023）系统梳理了ChatGPT在教育中的应用前景与风险，指出GAI在个性化学习、即时反馈和知识检索方面具有显著优势，但也存在信息准确性不足和学术诚信风险等问题。Mollick和Mollick（2023）提出了AI辅助教学的多种模式，包括AI作为导师、教练和模拟伙伴等角色，为教育实践提供了新的思路。国内学者刘红宇等（2024）调查发现，我国大学生AI工具使用率已超过60\%，但使用深度参差不齐，多数学生仍停留在简单的信息查询层面。

\subsubsection{AI使用行为与学习绩效研究}

关于AI使用行为与学习绩效的关系，现有研究呈现出两种不同的观点。持积极态度的学者认为，AI工具能够显著提升学习效率和学习质量。Becker等（2023）的实验研究表明，合理使用ChatGPT的学生在写作任务中的表现显著优于对照组。Dai等（2023）发现，AI辅助学习能够降低学生的认知负荷，使其将更多注意力投入到高阶思维活动中。持审慎态度的学者则指出，过度依赖AI可能产生负面效应。Darvishi等（2024）的研究发现，长期使用AI工具的学生在独立解题能力上出现明显下降，存在"认知卸载"现象。国内学者张伟和陈明（2024）的调查也表明，AI依赖程度较高的学生在创造性思维测试中的得分显著低于低依赖组。

\subsubsection{批判性思维与AI依赖研究}

批判性思维作为高等教育培养的核心能力之一，在AI时代被赋予了新的内涵。Pennycook和Rand（2020）指出，批判性思维是抵御AI生成错误信息的关键能力。Atkinson和Norris（2023）提出了"算法批判素养"的概念，强调学生需要具备评估AI输出质量的能力。AI依赖作为技术依赖的新形式，Tanaka等（2024）将其界定为"个体在面对学习任务时优先选择使用AI工具而非独立思考的倾向"，并发现AI依赖与自我效能感呈显著负相关。

\subsubsection{现有研究的特点与启示}

综合来看，现有研究存在以下不足：第一，多数研究将AI使用视为同质行为，忽视了使用深度差异对学习效果的不同影响；第二，缺乏将批判性思维与AI依赖同时纳入分析框架的研究，难以揭示二者在AI使用与学习绩效之间的复合作用机制；第三，研究对象多为西方高校学生，针对中国大学生尤其是经济学类专业学生的实证研究较为匮乏。本文正是基于上述研究缺口，构建"AI使用深度—学习绩效"的分析框架，引入批判性思维与AI依赖作为调节变量，以国际经济与贸易专业本科生为研究对象进行实证检验。

\subsection{研究的主要工作}

\subsubsection{理论梳理与框架构建}

系统回顾和评述技术接受理论、认知负荷理论与建构主义学习理论的核心观点，结合生成式人工智能的技术特征与教育应用场景，构建"AI使用深度—学习绩效"的理论分析框架，明确批判性思维与AI依赖的调节作用路径。

\subsubsection{研究设计与问卷开发}

采用实证研究方法，以问卷调查为主要数据收集手段。在文献研究基础上，借鉴成熟量表并结合本研究情境，开发AI使用深度、批判性思维、AI依赖和学习绩效的测量工具，并通过预调查对问卷进行修正与优化。

\subsubsection{数据收集与样本选择}

将目标调查对象设定为高校国际经济与贸易专业本科生，采用分层随机抽样方法，通过线上与线下相结合的方式发放问卷，确保样本的代表性与数据质量。

\subsubsection{数据分析与假设检验}

运用SPSS与Stata等专业统计分析软件对收集到的一手数据进行系统分析。首先进行描述性统计与信效度检验，确保数据质量；其次通过相关分析初步验证变量间关系；最后运用多元回归分析与调节效应检验，验证研究假设。

\subsubsection{研究结论与展望}

依据数据分析结果归纳得出研究结论，提出针对性的教育政策建议，并讨论本研究的局限性及未来研究方向。

\section{理论背景与研究假设}

\subsection{AI使用深度的概念界定}

AI使用深度是本文的核心概念，其与传统"使用频率"有着本质区别。使用频率仅反映学生使用AI工具的次数与时长，而AI使用深度强调的是学生在使用过程中的认知参与程度、互动质量以及高阶学习投入水平。借鉴Heerink等（2010）关于技术接受深度的研究以及Webster和Hackley（1997）关于学习投入的理论，本文将AI使用深度界定为三个维度：一是认知参与深度，即学生在使用AI时是否进行主动思考与批判性评估，而非被动接受AI输出；二是互动质量深度，即学生与AI的交互是否具有多轮对话、追问探究和反馈修正等特征；三是高阶学习投入深度，即学生是否将AI用于分析、综合、评价等高阶认知活动，而非仅用于信息检索和简单问答。

一个关键且有多维度特性的概念，AI使用深度的提出有助于超越"使用与否"的二元框架，从过程视角揭示AI技术影响学习绩效的内在机制。

\subsection{AI使用深度与学习绩效的关系}

认知负荷理论认为，学习过程中的认知资源是有限的，合理分配认知资源是提升学习绩效的关键。当学生以深度方式使用AI时，AI承担了信息检索和初步加工等低阶认知任务，使学生能够将有限的认知资源集中于理解、分析和评价等高阶认知活动，从而优化认知负荷分配，提升学习绩效。建构主义学习理论也支持这一观点，深层次的AI使用本质上是一种建构性学习过程，学生通过与AI的深度互动，不断修正和完善自己的知识结构。

据此，本研究提出以下假设：

H1：AI使用深度对大学生学习绩效具有显著正向影响。

H1a：认知参与深度对大学生学习绩效具有显著正向影响。

H1b：互动质量深度对大学生学习绩效具有显著正向影响。

H1c：高阶学习投入深度对大学生学习绩效具有显著正向影响。

\subsection{批判性思维的调节作用}

批判性思维是指个体在面对信息时进行理性分析、逻辑推理和独立判断的能力。在AI使用情境中，批判性思维水平较高的学生更倾向于对AI生成的内容进行质疑和验证，而非盲目接受。这种批判性审视有助于学生从AI辅助中获取更高质量的知识，同时避免被AI的错误输出所误导。相反，批判性思维水平较低的学生可能将AI输出视为权威答案，导致浅层学习甚至认知错误。

过往研究已表明，批判性思维在技术使用与学习效果之间发挥重要的调节作用。在AI使用深度与学习绩效的关系中，批判性思维可能起到"放大器"的作用：当学生具备较高水平的批判性思维时，AI使用深度对学习绩效的正向影响更为显著；当批判性思维水平较低时，即使AI使用深度较高，其对学习绩效的促进作用也可能受限。

基于上述讨论，本研究提出以下假设：

H2：批判性思维在AI使用深度与学习绩效之间起正向调节作用，即批判性思维水平越高，AI使用深度对学习绩效的正向影响越强。

\subsection{AI依赖的调节作用}

AI依赖是指个体在面对学习任务时优先选择使用AI工具而非独立思考的倾向。与AI使用深度不同，AI依赖强调的是一种心理状态和行为习惯，而非使用过程中的认知投入水平。高AI依赖的学生倾向于将AI视为"答案提供者"，而非"学习伙伴"，即使在使用过程中也可能缺乏深度思考。

AI依赖可能在AI使用深度与学习绩效之间起负向调节作用。当AI依赖程度较高时，学生可能将AI使用深度中的认知参与和互动探究转化为对AI的进一步依赖，反而削弱了独立思考能力，降低学习绩效。当AI依赖程度较低时，学生能够更好地控制AI使用的边界，将深度使用转化为真正的学习收益。

H3：AI依赖在AI使用深度与学习绩效之间起负向调节作用，即AI依赖程度越高，AI使用深度对学习绩效的正向影响越弱。

\section{研究方法}

\subsection{调查设计与数据收集}

为了保证调查问卷的信度与效度，本研究采用以下程序进行问卷设计与数据收集。首先，在文献研究基础上，借鉴国内外成熟量表，结合国际经济与贸易专业本科生的学习情境，初步编制调查问卷。其次，邀请5名教育学与经济学领域的专家对问卷内容进行审阅，并根据专家意见进行修改。再次，通过预调查（发放问卷50份）检验问卷的表述清晰度和内部一致性，对部分题项进行调整。最终形成正式问卷，包含AI使用深度量表（12题项）、批判性思维量表（7题项）、AI依赖量表（6题项）和学习绩效量表（8题项），共33个题项，均采用Likert五级量表（1=完全不同意，5=完全同意）。

本研究以南京、上海、杭州三地高校国际经济与贸易专业本科生为调查对象，采用分层随机抽样方法，综合考虑学校层次（985/211/普通本科）和年级（大一至大四）进行分层。通过线上（问卷星平台）与线下（课堂发放）相结合的方式发放问卷，共发放问卷450份，回收有效问卷386份，有效回收率为85.8\%。

\subsection{变量的选取与模型设定}

\subsubsection{被解释变量}

学习绩效（LP）：借鉴Hase和Kenyon（2007）的学习绩效量表，从知识获取、技能提升和学习满意度三个维度进行测量，共8个题项。Cronbach's $\alpha$系数为0.892，表明量表具有较好的内部一致性。

\subsubsection{解释变量}

AI使用深度（AUD）：本文自行开发的量表，包含认知参与深度（4题项）、互动质量深度（4题项）和高阶学习投入深度（4题项）三个维度，共12个题项。Cronbach's $\alpha$系数为0.913，各维度的$\alpha$系数分别为0.876、0.854和0.891。

\subsubsection{调节变量}

批判性思维（CT）：借鉴Facione（1990）的California Critical Thinking Disposition Inventory（CCTDI）量表，选取7个题项进行测量。Cronbach's $\alpha$系数为0.845。

AI依赖（AID）：借鉴Tanaka等（2024）的技术依赖量表，结合AI使用情境进行改编，共6个题项。Cronbach's $\alpha$系数为0.831。

\subsubsection{控制变量}

本研究控制了以下变量：性别（Gender，男=1，女=0）、年级（Grade，大一至大四分别赋值1--4）、GPA（学业成绩，1--5级）、每日AI使用时长（Time，1=不足30分钟，2=30--60分钟，3=1--2小时，4=2小时以上）。

相关变量及赋值见表3-1。

\begin{table}[htbp]
  \centering
  \caption{模型中相关变量及其赋值情况}
  \begin{tabular}{llll}
    \toprule
    变量类型 & 变量名称 & 变量符号 & 赋值说明 \\
    \midrule
    被解释变量 & 学习绩效 & LP & Likert五级量表均值 \\
    \multirow{3}{*}{解释变量} & 认知参与深度 & AUD1 & Likert五级量表均值 \\
     & 互动质量深度 & AUD2 & Likert五级量表均值 \\
     & 高阶学习投入深度 & AUD3 & Likert五级量表均值 \\
    \multirow{2}{*}{调节变量} & 批判性思维 & CT & Likert五级量表均值 \\
     & AI依赖 & AID & Likert五级量表均值 \\
    \multirow{4}{*}{控制变量} & 性别 & Gender & 男=1，女=0 \\
     & 年级 & Grade & 大一至大四赋值1--4 \\
     & 学业成绩 & GPA & 1--5级 \\
     & 每日AI使用时长 & Time & 1--4级 \\
    \bottomrule
  \end{tabular}
\end{table}

\subsubsection{模型设定}

为检验研究假设，本文构建以下计量模型：

基准回归模型：
$$LP_i = \beta_0 + \beta_1 AUD_i + \beta_2 Gender_i + \beta_3 Grade_i + \beta_4 GPA_i + \beta_5 Time_i + \varepsilon_i$$

调节效应模型：
$$LP_i = \beta_0 + \beta_1 AUD_i + \beta_2 CT_i + \beta_3 AUD_i \times CT_i + \beta_4 Gender_i + \beta_5 Grade_i + \beta_6 GPA_i + \beta_7 Time_i + \varepsilon_i$$

$$LP_i = \beta_0 + \beta_1 AUD_i + \beta_2 AID_i + \beta_3 AUD_i \times AID_i + \beta_4 Gender_i + \beta_5 Grade_i + \beta_6 GPA_i + \beta_7 Time_i + \varepsilon_i$$

其中，$LP_i$为第$i$个学生的学习绩效，$AUD_i$为AI使用深度，$CT_i$为批判性思维，$AID_i$为AI依赖，$AUD_i \times CT_i$和$AUD_i \times AID_i$分别为交互项，$\varepsilon_i$为随机误差项。

\section{研究分析与发现}

\subsection{描述性统计分析}

根据调查数据，表4-1报告了样本数据的人口统计特征。在386个有效样本中，男生占42.5\%，女生占57.5\%；大一至大四分别占22.3\%、25.4\%、28.2\%和24.1\%；GPA分布较为均匀，3分以上占比超过70\%。每日AI使用时长方面，使用1--2小时的学生最多，占38.9\%，使用2小时以上的占21.5\%，表明国际经济与贸易专业本科生AI使用较为普遍。

\begin{table}[htbp]
  \centering
  \caption{样本数据人口统计特征}
  \begin{tabular}{llcc}
    \toprule
    统计特征 & 类别 & 频数 & 百分比（\%） \\
    \midrule
    \multirow{2}{*}{性别} & 男 & 164 & 42.5 \\
     & 女 & 222 & 57.5 \\
    \midrule
    \multirow{4}{*}{年级} & 大一 & 86 & 22.3 \\
     & 大二 & 98 & 25.4 \\
     & 大三 & 109 & 28.2 \\
     & 大四 & 93 & 24.1 \\
    \midrule
    \multirow{4}{*}{每日AI使用时长} & 不足30分钟 & 56 & 14.5 \\
     & 30--60分钟 & 97 & 25.1 \\
     & 1--2小时 & 150 & 38.9 \\
     & 2小时以上 & 83 & 21.5 \\
    \bottomrule
  \end{tabular}
\end{table}

表4-2报告了主要变量的描述性统计结果。AI使用深度均值为3.42（标准差0.71），其中认知参与深度均值3.38、互动质量深度均值3.35、高阶学习投入深度均值3.52，表明样本学生的AI使用深度总体处于中等偏上水平。学习绩效均值为3.56（标准差0.68），批判性思维均值为3.29（标准差0.74），AI依赖均值为2.87（标准差0.82），说明样本学生AI依赖程度相对较低。

\begin{table}[htbp]
  \centering
  \caption{主要变量描述性统计}
  \begin{tabular}{lcccc}
    \toprule
    变量 & 均值 & 标准差 & 最小值 & 最大值 \\
    \midrule
    学习绩效（LP） & 3.56 & 0.68 & 1.25 & 5.00 \\
    AI使用深度（AUD） & 3.42 & 0.71 & 1.17 & 5.00 \\
    认知参与深度（AUD1） & 3.38 & 0.76 & 1.00 & 5.00 \\
    互动质量深度（AUD2） & 3.35 & 0.73 & 1.00 & 5.00 \\
    高阶学习投入深度（AUD3） & 3.52 & 0.74 & 1.25 & 5.00 \\
    批判性思维（CT） & 3.29 & 0.74 & 1.00 & 5.00 \\
    AI依赖（AID） & 2.87 & 0.82 & 1.00 & 4.83 \\
    \bottomrule
  \end{tabular}
\end{table}

\subsection{信度与效度检验}

本研究采用Cronbach's $\alpha$系数检验量表的内部一致性信度。结果显示，AI使用深度量表的$\alpha$系数为0.913，其三个子维度的$\alpha$系数分别为0.876、0.854和0.891；批判性思维量表的$\alpha$系数为0.845；AI依赖量表的$\alpha$系数为0.831；学习绩效量表的$\alpha$系数为0.892。所有量表的$\alpha$系数均大于0.8，表明量表具有较好的内部一致性。

效度检验方面，本研究采用验证性因子分析（CFA）检验量表的建构效度。结果显示，AI使用深度三因子模型的拟合指标良好（$\chi^2/df$=2.34，RMSEA=0.059，CFI=0.946，TLI=0.935，SRMR=0.042），各题项的标准化因子载荷均大于0.6且在$p<0.001$水平上显著，平均方差提取量（AVE）均大于0.5，组合信度（CR）均大于0.8，表明量表具有较好的收敛效度。各构念AVE的平方根均大于其与其他构念的相关系数，表明区分效度良好。

\subsection{相关分析}

Pearson相关分析结果显示，AI使用深度与学习绩效呈显著正相关（$r$=0.463，$p<0.01$），批判性思维与学习绩效呈显著正相关（$r$=0.387，$p<0.01$），AI依赖与学习绩效呈显著负相关（$r$=-0.215，$p<0.01$），AI使用深度与AI依赖呈弱正相关（$r$=0.128，$p<0.05$）。相关分析结果初步支持了研究假设，为后续回归分析奠定了基础。

\subsection{回归分析结果}

这一部分报告多元回归分析结果。为避免多重共线性问题，在构建交互项之前对AI使用深度、批判性思维和AI依赖进行了中心化处理。各模型的方差膨胀因子（VIF）均小于5，不存在严重的多重共线性问题。

\subsubsection{AI使用深度对学习绩效的主效应检验}

表4-3报告了AI使用深度对学习绩效的基准回归结果。模型1仅包含控制变量，模型2加入AI使用深度，模型3将AI使用深度拆分为三个子维度。

\begin{table}[htbp]
  \centering
  \caption{AI使用深度对学习绩效的回归结果}
  \begin{tabular}{lccc}
    \toprule
    变量 & 模型1 & 模型2 & 模型3 \\
    \midrule
    性别（Gender） & 0.052 & 0.038 & 0.041 \\
    年级（Grade） & 0.067 & 0.043 & 0.039 \\
    GPA & 0.284*** & 0.213*** & 0.198*** \\
    每日AI使用时长（Time） & 0.091* & 0.047 & 0.042 \\
    AI使用深度（AUD） &  & 0.412*** &  \\
    认知参与深度（AUD1） &  &  & 0.156** \\
    互动质量深度（AUD2） &  &  & 0.134** \\
    高阶学习投入深度（AUD3） &  &  & 0.187*** \\
    \midrule
    $R^2$ & 0.112 & 0.276 & 0.298 \\
    调整$R^2$ & 0.103 & 0.267 & 0.285 \\
    $F$值 & 12.07*** & 29.14*** & 22.86*** \\
    \bottomrule
    \multicolumn{4}{l}{\footnotesize 注：*表示$p<0.1$，**表示$p<0.05$，***表示$p<0.01$}
  \end{tabular}
\end{table}

模型2结果显示，AI使用深度对学习绩效具有显著正向影响（$\beta$=0.412，$p<0.01$），假设H1得到支持。模型3进一步显示，高阶学习投入深度的影响最大（$\beta$=0.187，$p<0.01$），认知参与深度次之（$\beta$=0.156，$p<0.05$），互动质量深度也有显著正向影响（$\beta$=0.134，$p<0.05$），假设H1a、H1b、H1c均得到支持。

\subsubsection{批判性思维的调节效应检验}

表4-4报告了批判性思维在AI使用深度与学习绩效之间的调节效应检验结果。

\begin{table}[htbp]
  \centering
  \caption{批判性思维的调节效应检验结果}
  \begin{tabular}{lcc}
    \toprule
    变量 & 模型4 & 模型5 \\
    \midrule
    AI使用深度（AUD） & 0.356*** & 0.341*** \\
    批判性思维（CT） & 0.224*** & 0.218*** \\
    AUD$\times$CT &  & 0.156** \\
    控制变量 & 是 & 是 \\
    \midrule
    $R^2$ & 0.342 & 0.365 \\
    $\Delta R^2$ &  & 0.023** \\
    $F$值 & 28.73*** & 27.46*** \\
    \bottomrule
    \multicolumn{3}{l}{\footnotesize 注：**表示$p<0.05$，***表示$p<0.01$}
  \end{tabular}
\end{table}

模型5结果显示，交互项AUD$\times$CT的系数显著为正（$\beta$=0.156，$p<0.05$），$\Delta R^2$=0.023且在$p<0.05$水平上显著，表明批判性思维在AI使用深度与学习绩效之间起显著正向调节作用，假设H2得到支持。这意味着，对于批判性思维水平较高的学生，AI使用深度对学习绩效的促进作用更为显著。

\subsubsection{AI依赖的调节效应检验}

表4-5报告了AI依赖在AI使用深度与学习绩效之间的调节效应检验结果。

\begin{table}[htbp]
  \centering
  \caption{AI依赖的调节效应检验结果}
  \begin{tabular}{lcc}
    \toprule
    变量 & 模型6 & 模型7 \\
    \midrule
    AI使用深度（AUD） & 0.398*** & 0.412*** \\
    AI依赖（AID） & -0.176*** & -0.168*** \\
    AUD$\times$AID &  & -0.142** \\
    控制变量 & 是 & 是 \\
    \midrule
    $R^2$ & 0.318 & 0.337 \\
    $\Delta R^2$ &  & 0.019** \\
    $F$值 & 26.54*** & 24.89*** \\
    \bottomrule
    \multicolumn{3}{l}{\footnotesize 注：**表示$p<0.05$，***表示$p<0.01$}
  \end{tabular}
\end{table}

模型7结果显示，交互项AUD$\times$AID的系数显著为负（$\beta$=-0.142，$p<0.05$），$\Delta R^2$=0.019且在$p<0.05$水平上显著，表明AI依赖在AI使用深度与学习绩效之间起显著负向调节作用，假设H3得到支持。这意味着，AI依赖程度越高，AI使用深度对学习绩效的正向影响越弱。

\subsection{稳健性检验}

为确保研究结论的可靠性，本文进行了以下稳健性检验：第一，替换被解释变量，将学习绩效的三个子维度分别作为被解释变量进行回归，核心结论保持一致；第二，采用分样本回归，按年级和学校层次分别进行回归分析，AI使用深度的正向影响在各子样本中均显著；第三，采用Bootstrap方法检验调节效应（重复抽样5000次），批判性思维与AI依赖的调节效应均通过显著性检验。上述稳健性检验结果表明，本文的核心研究结论具有较强的稳健性。

\section{研究结论与展望}

\subsection{研究结论}

本研究通过构建"AI使用深度—学习绩效"的理论分析框架，以高校国际经济与贸易专业本科生为研究对象，运用问卷调查法与多元回归分析法，系统考察了AI使用深度对学习绩效的影响以及批判性思维与AI依赖的调节作用，得出以下主要结论：

第一，AI使用深度对大学生学习绩效具有显著正向影响。这一发现表明，AI对学习的影响不应仅从"是否使用"或"使用频率"的角度来评估，更应关注使用的"深度"与"质量"。深层次、互动式、反思型的AI使用能够有效促进知识理解与学习能力提升，其中高阶学习投入深度的促进作用最为显著。

第二，批判性思维在AI使用深度与学习绩效之间起显著正向调节作用。批判性思维水平较高的学生能够更好地利用AI工具进行深度学习，对AI输出进行批判性审视和创造性转化，从而获得更大的学习收益。这一发现凸显了在AI时代培养大学生批判性思维的重要性。

第三，AI依赖在AI使用深度与学习绩效之间起显著负向调节作用。AI依赖程度较高的学生，即使在使用深度较高的情况下，其学习绩效的提升也受到限制。这说明AI依赖可能削弱学生的独立思考能力和自主性，从而抵消深度使用带来的部分学习收益。

\subsection{政策建议}

基于上述研究结论，本文提出以下政策建议：

第一，高校应加强大学生AI素养教育。将AI素养纳入通识教育课程体系，培养学生对AI工具的批判性使用能力，使学生既能有效利用AI提升学习效率，又能保持独立思考与判断能力。

第二，优化课程设计与教学评价体系。在课程设计中增加需要深度思考和创造性解决问题的任务，减少可以通过AI简单完成的机械性作业。在教学评价中，注重过程性评价与多元评价，降低AI代写等学术不端行为的激励。

第三，引导学生形成"人机协同"的学习模式。鼓励学生将AI视为"学习伙伴"而非"答案机器"，通过多轮对话、追问探究和反思修正等方式与AI进行深度互动，实现人机优势互补。

\subsection{局限性与展望}

本研究存在以下局限性：第一，采用横截面数据，难以揭示AI使用深度与学习绩效之间的因果关系，未来研究可采用纵向追踪设计；第二，研究对象仅限于国际经济与贸易专业本科生，研究结论的推广性有待进一步验证；第三，AI使用深度的测量依赖自我报告，可能存在社会期望偏差，未来研究可结合日志数据等客观指标进行补充测量。此外，随着GAI技术的持续演进，AI使用深度和学习绩效的内涵可能发生变化，后续研究应关注技术发展对理论框架的动态影响。

\makereference

\begin{enumerate}[label={[\arabic*]}]
  \item Kasneci E, Sessler K, K\"uchemann S, et al. ChatGPT for good? On opportunities and challenges of large language models for education[J]. Learning and Individual Differences, 2023, 103: 102274.
  \item Mollick E R, Mollick L. Using AI to implement effective teaching strategies in classrooms: Five strategies, including prompts[J]. The Wharton School Research Paper, 2023.
  \item Becker B, Dohmen P, Engel J, et al. ChatGPT and the end of homework: Rethinking student assessment in the age of AI[J]. Educational Technology Research and Development, 2023, 71(4): 1023-1041.
  \item Dai Y, Chai C S, Lin P Y, et al. AI-powered educational chatbots for facilitating student learning: A systematic review[J]. Computers and Education, 2023, 198: 104758.
  \item Darvishi A, Khosravi H, Sadiq S, et al. Impact of AI assistance on student agency[J]. Computers and Education, 2024, 210: 104967.
  \item Heerink M, Kr\"ose B, Evers V, et al. Assessing acceptance of assistive social agent technology by older adults: The Almere model[J]. International Journal of Social Robotics, 2010, 2(4): 361-375.
  \item Webster J, Hackley P. Teaching effectiveness in technology-mediated distance learning[J]. Academy of Management Journal, 1997, 40(6): 1282-1309.
  \item Hase S, Kenyon C. From andragogy to heutagogy[J]. UltiBase Articles, 2007, 5(3): 1-10.
  \item Facione P A. The California Critical Thinking Disposition Inventory[M]. Millbrae, CA: California Academic Press, 1990.
  \item Pennycook G, Rand D G. Lazy, not biased: Susceptibility to partisan fake news is better explained by lack of reasoning than by motivated reasoning[J]. Cognition, 2020, 188: 39-50.
  \item Atkinson L, Norris M. Algorithmic critical literacy: A framework for understanding and evaluating algorithms[J]. Journal of Media Literacy Education, 2023, 15(2): 45-62.
  \item Tanaka K, Yamamoto S, Ito T. AI dependency in higher education: Conceptualization and measurement[J]. Computers in Human Behavior, 2024, 152: 108076.
  \item 刘红宇, 张明华, 王丽丽. 大学生生成式人工智能使用行为调查与影响因素分析[J]. 高等教育研究, 2024, 45(3): 67-78.
  \item 张伟, 陈明. 生成式人工智能依赖对大学生创造性思维的影响研究[J]. 教育科学, 2024, 40(1): 45-53.
  \item 中国互联网络信息中心. 第53次中国互联网络发展状况统计报告[R]. 北京: 中国互联网络信息中心, 2024.
\end{enumerate}

\makeacknowledgement

论文顺利完成之际，我要衷心感谢我的导师李四教授。从论文选题、框架构建到数据分析与论文修改，李老师都给予了悉心指导和耐心帮助，其严谨的治学态度和渊博的学识使我受益匪浅。同时，感谢国际经济与贸易专业的各位老师，是你们的悉心教导为我打下了坚实的专业基础。感谢参与问卷调查的各位同学，你们的支持是本研究得以完成的重要保障。最后，感谢我的家人和朋友，你们的理解与鼓励是我不断前行的动力源泉。

\end{spacing}

\end{document}
